\section{The Soul and Its Work}

In Kabbalah the soul is not one thing. It is a \textbf{layered structure} of five
levels, each rooted in a different world, and it can migrate across lifetimes. The
human is not a passive creature in the cosmic account. The human is the
\textbf{active repairer}. Through right deeds the soul lifts the fallen sparks and
mends the broken world.

\subsection{The Five Levels of Soul}

The full count is \textbf{five levels}, often abbreviated by their initials as
NaRaN-ChaY. Each is rooted in one of the worlds, and each carries a different
faculty.

\begin{longtable}{L{0.16\linewidth} L{0.32\linewidth} L{0.20\linewidth} L{0.16\linewidth}}
\caption{The five levels of the soul and the world each is rooted in.}\\
\toprule
\textbf{Level} & \textbf{Plain meaning} & \textbf{Faculty} & \textbf{Rooted in} \\
\midrule
Nefesh & rest, the animal soul & life-force, instinct, the vitality of the body & Assiah \\
Ruach & spirit, wind & emotion, moral character & Yetzirah \\
Neshamah & breath, soul & intellect, divine consciousness & Beriah \\
Chayah & the living one & the bridge to the divine will & Atzilut \\
Yechidah & the singular one & the spark of pure unity with God & Adam Kadmon \\
\bottomrule
\end{longtable}

A person normally accesses only the \textbf{lower three}, nefesh, ruach, and
neshamah. The upper two, \textbf{chayah and yechidah}, are usually transcendent
surrounding lights, not drawn fully inside. They are earned only by great
spiritual work. The soul reaches up through its levels as the cosmos reaches down
through its worlds.

\subsection{Gilgul and Ibbur}

The soul migrates and combines. Two named mechanisms run the traffic of souls
across lives.

\begin{itemize}
 \item \textbf{Reincarnation} (gilgul, cycle or rolling). A soul returns in a new
 body to complete the rectification it failed to finish, the commandments
 it left undone. Each return is a fresh chance at the unfinished work.
 \item \textbf{Impregnation} (ibbur). A second soul joins a living person for a
 time, to help complete a task or gain a benefit. It is a positive joining.
 The added soul leaves when its work is done.
\end{itemize}

Every individual soul belongs to a larger \textbf{soul-root}, and the root
subdivides into many \textbf{sparks}. Souls that share a root are linked across
generations. An ibbur often draws a soul from the same root.

\subsection{The 288 Sparks}

When the \textbf{vessels shattered} in the creation account, the divine light
scattered. The standard count of fallen sparks is \textbf{288}. These sparks are
trapped within the shells (qelippot) and within the material world. The task of
humanity is the \textbf{sifting} (birur), the gathering and elevating of those fallen
sparks back to their source. Each commandment done with intention raises sparks.
This is the engine of repair.

\subsection{The Body as the Image of the Divine}

The human body is read as the \textbf{image of the divine}, the form of the
primordial Man and the shape of the sefirot. \enquote{God created man in His
image} is read as a structural claim. The body mirrors the tree.

The \textbf{613 commandments} are mapped onto the body. The total splits into two
classes that match the two kinds of body-part.

\begin{longtable}{L{0.26\linewidth} L{0.12\linewidth} L{0.44\linewidth}}
\caption{The 613 commandments mapped to the limbs and sinews.}\\
\toprule
\textbf{Class} & \textbf{Count} & \textbf{Mapped to} \\
\midrule
Positive commandments & 248 & the 248 limbs and bones of the body \\
Negative commandments & 365 & the 365 sinews of the body, and the 365 days of the solar year \\
Total & 613 & the whole body \\
\bottomrule
\end{longtable}

The mapping is exact. Every positive commandment animates a \textbf{limb}. Every
negative commandment guards a \textbf{sinew}. Since 248 plus 365 equals 613, the
whole correspondence turns on a single clean identity. The body is a map of the
law, and the law is a map of the divine form.

\subsection{The Human as the Active Repairer}

This is the center of the Kabbalistic account of the human. The soul is planted in a
broken world for a reason. Its deeds can find and free the fallen sparks. This
work is the \textbf{repair of the world} (tikkun olam), in its original cosmic
sense, the gathering of the scattered light and the mending of the broken vessels.

Two practices drive the repair.

\begin{itemize}
 \item \textbf{Intention} (kavanah). The inner aim that a deed is performed
 with. A commandment done with the right intention raises a spark. The same
 deed done emptily does not. Intention is what gives an action its
 spark-lifting power.
 \item \textbf{Unification} (yichud). The drawing-together of the divided
 aspects of the divine, above and below, the higher source and the exiled
 presence. The aim of right action is to heal the cosmic split.
\end{itemize}

So the soul is not a spectator. It is the worker on whom the whole repair depends.
Creation broke. The mending is assigned to the human being, and the tools of the
mending are intention and unification. Every well-aimed deed is a small reversal of
the scattering, a single spark carried home.
